%-------------------------------------------------------------------------------
\begin{abstract}
Benchmark and system parameters often have a significant impact
on performance evaluation, which raises
a long-lasting question about which settings we should use.

This paper studies the feasibility and benefits of extensive
evaluation, by reproducing 10 prior works and testing them under
a wider range of settings. We make two contributions: first, we
show that sampling and regression is effective to predict performance
for well-engineered systems. Based on this observation, to reduce the cost
of extensive evaluation, we suggest gradually testing more settings until the
prediction accuracy becomes stable. Second, we design and develop a tool to
visualize high-dimensional data from extensive evaluation on a figure called Tiled Heat Map.

Our experience shows that extensive evaluation can achieve two goals:
first, it has revealed many correctness and performance issues, and
addressing them can improve the quality of artifact and facilitate future
comparison. Second, extensive evaluation can generate a more comprehensive
comparison among different works.


\end{abstract}